% !TeX root = ../technical_paper.tex

\section{总结与展望}

\subsection{工作总结}

本文围绕多 MCU 协同的光储一体化数字电源系统开展了设计、实现与测试工作。系统以 GD32 系列 MCU 为核心控制平台，构建了由主控制器、BMS 控制器和无线通信控制器组成的分布式控制架构。其中，主控制板采用 GD32G553，负责系统状态调度、模拟量采样、控制接口输出、通信数据汇总和故障状态管理；BMS 板采用 GD32C113，结合 GD30BM2016 电池监测芯片完成 16S 磷酸铁锂电池组的电压、电流、温度采集、保护判断、均衡控制和电池状态上报；无线监控模块采用 GD32VW553，负责 WiFi 通信、上位机交互和 App 状态显示。通过不同 GD32 MCU 的协同工作，系统实现了主控、储能安全管理和无线监控功能的分离，提高了平台的可维护性、可扩展性和工程实现灵活性。

基于 GD32 MCU 构建整套控制平台，使系统在控制实时性、外设资源利用和国产化器件应用方面具有较好的工程意义。GD32G553 提供了丰富的 ADC、HRTIMER、CAN、UART 和 GPIO 资源，能够满足数字电源系统对多通道采样、PWM 控制、故障输入和高速通信的需求；GD32C113 具备较好的低成本控制能力和通信接口资源，适合承担 BMS 状态机、保护逻辑和数据遥测任务；GD32VW553 集成无线通信能力，能够将传统嵌入式控制系统扩展到上位机和移动端监控场景。三类 GD32 MCU 分别承担不同任务，使系统既保留了嵌入式实时控制能力，又具备远程可视化和智能交互能力。

在硬件方面，本文完成了主控制板、BMS 板和无线监控模块的原理设计与功能验证。主控制板完成了 GD32G553 最小系统、ADC 采样接口、PWM 输出接口、硬件故障输入、CAN 通信接口和 UART 通信接口设计；BMS 板完成了 GD32C113 主控电路、GD30BM2016 AFE 采样电路、电流检测电路、温度采样电路、均衡电路、充放电 MOSFET 驱动与保护电路设计；无线模块完成了与主控制器的数据交互接口设计。辅助电源部分实现了电池包输入到 12V、5V 和 3.3V 的多级供电，为控制板和通信模块提供稳定电源。功率器件选型围绕 16S 磷酸铁锂电池组的电压范围进行，兼顾耐压裕量、导通损耗、热设计和保护可靠性。

在软件方面，本文完成了基于 GD32 MCU 的多节点协同软件设计。主控制器软件采用分层架构，包含外设驱动、采样处理、通信协议、状态机管理和控制接口管理等模块；BMS 软件围绕产品状态机运行，完成 AFE 初始化、周期测量、保护判断、SOC 估算、均衡控制、FET 策略和 CAN 遥测上报；无线监控软件实现了主控数据接收、状态上传、上位机控制和 App 显示功能。系统能够通过上位机或 App 查看电池电压、电流、SOC、温度、工作模式和故障状态，并能够下发模式切换和参数配置命令。该软件架构体现了 GD32 MCU 在多节点嵌入式协同控制中的应用价值。

在测试方面，本文完成了硬件基础功能、采样与保护、通信链路、上位机与 App 功能、多模式切换和低功率充放电实物测试。测试结果表明，主控制板、BMS 板和无线模块能够正常上电运行，CAN、UART 和 WiFi 通信链路能够完成数据交互，BMS 数据能够经主控制器上传至上位机和 App，系统能够根据 BMS 允许状态完成待机、运行、充电、放电、降级和故障等模式切换。当 BMS 上报禁止充电、禁止放电或故障状态时，主控制器能够限制相关动作并通过显示端提示异常，体现了安全优先的设计原则。

本文工作的意义在于，基于 GD32 MCU 构建了一个面向光储一体化电源系统的国产化多控制器协同平台。该平台不仅完成了电池管理、无线监控、上位机交互和多模式状态切换，也为后续主功率级闭环控制和整机工程化实现提供了可复用的软件和硬件基础。相比单控制器集中式方案，本系统通过 GD32G553、GD32C113 和 GD32VW553 的功能分工降低了系统复杂度，使各模块能够独立调试、独立验证并逐步扩展，符合复杂电源系统分层设计的发展方向。

需要说明的是，本文当前阶段重点完成了控制、采样、保护、通信和监控平台的搭建与验证，三端口功率变换主回路尚未完成完整实物闭环实现。因此，本文未对三端口主功率级的满功率效率、额定负载温升和长时间闭环运行性能作结论性评价。现有工作为后续主功率级接入、闭环控制调试和整机性能测试奠定了软硬件基础。

\subsection{不足与改进方向}

虽然系统已经完成阶段性设计和验证，但仍存在需要进一步完善的内容。

首先，当前测试主要围绕控制板、BMS 板、无线通信模块、上位机交互、App 显示和低功率充放电等内容展开，系统在更高功率等级、更复杂负载条件和更长时间连续运行条件下的性能仍需进一步验证。后续可在现有控制、采样、保护和通信平台基础上，逐步开展更完整的功率级接入、闭环控制调试、额定工况测试和长时间稳定性测试，从而进一步评价系统在实际应用场景下的效率、温升、纹波和动态响应能力。

其次，当前充放电测试仍处于低功率验证阶段，主要用于验证主控、BMS、无线监控和保护逻辑之间的协同关系。后续需要在安全条件下逐步提升测试功率等级，增加不同电池 SOC、不同负载功率、不同温度条件和不同工作模式下的测试样本，并记录系统效率、输出稳定性、负载阶跃响应和关键器件温升等数据，为后续系统优化提供更充分的实验依据。

再次，无线监控和 App 显示功能已经具备基础状态展示和模式控制能力，但界面交互、历史数据记录、故障日志管理和参数可视化配置仍可进一步优化。后续可以增加运行数据存储、故障记录导出、曲线显示、远程参数整定和异常事件回放等功能，使系统不仅能够实时显示当前状态，也能够对历史运行过程进行分析，从而提升系统调试、维护和演示效果。

最后，系统保护策略还需要在更多异常场景下进行验证。例如 BMS 通信中断、单体电压突变、温度传感器异常、无线链路掉线、主控复位、驱动异常和采样信号漂移等情况，都需要进一步完善测试用例。后续可在现有软硬件平台基础上建立更完整的故障注入测试流程，通过软硬件联合保护提高系统在复杂工况下的安全性和可靠性。

总体来看，现有工作已经完成了控制、采样、保护、通信和监控平台的基础搭建，为后续更高功率等级的系统验证、复杂能量调度策略实现和多端口功率变换方案拓展提供了良好基础。
\subsection{后续工作展望}

后续工作将围绕主功率级实现、闭环控制完善、系统级测试和人机交互优化展开。首先，将在现有 GD32G553 主控制板和 GD32C113 BMS 平台基础上完成主功率级硬件接入，逐步开展低压限流上电、PWM 驱动验证、采样校准和保护阈值测试。其次，将结合实测波形调整控制参数，完善充电、放电和模式切换过程中的闭环控制策略，提高系统动态响应和稳定性。再次，将开展额定功率条件下的效率、温升、纹波和负载阶跃测试，形成更完整的性能评价数据。最后，将继续完善基于 GD32VW553 的无线通信和 App 功能，实现更友好的状态显示、故障分析和参数配置。

总体而言，本文完成了基于 GD32 MCU 的光储一体化数字电源系统控制、采样、保护、通信和监控基础平台搭建，验证了 GD32 多 MCU 协同架构在储能电源系统中的可行性。虽然主功率级仍需进一步完善，但现有成果已经为后续整机闭环实现和工程化优化提供了可靠基础。